\section{Risks and limitations}
\label{sec:risks}

This section assesses the principal failure modes. Six are
structural; none is fatal; two reorder the roadmap.

\subsection{The trusted computing base, itemized}
\label{sec:risks-tcb}

The design of \cref{sec:dsl} moves almost everything out of the TCB;
what remains warrants explicit enumeration. Trusted for
\emph{soundness of the theorems}: the Lean kernel; the goal
elaborator in \code{VesCore} (vocabulary term $\to$ Lean
\code{Prop} --- kept small, data-driven from the vocabulary source,
and cross-checked by recomputing statement hashes on both sides); and
the axiom packages (\code{RustAssumptions}, the trust-grant registry,
and the standard's own undefined-behavior latitude). Explicitly
\emph{untrusted}: the proc macro, the surface tactics (Lean
metaprograms whose output the kernel checks), all proof automation,
and \code{ves-check} itself --- a structural-rule engine whose failure
modes are false alarms or missed bookkeeping, not false theorems.

Trusted for \emph{correspondence} --- the claim that the theorem is
about the code in the repository --- are two mechanically narrowed
items. First, the extraction path: the macro and the checker link the
same \code{ves-syntax} parser, so there is one grammar, and the
subject hash covers the protected expression's exact token stream.
Second, the token-identity property: the macro must re-emit the
\code{then}-expression verbatim, which is enforced by expansion
snapshots in \code{ves-macros}' own CI. What no hash can close is the
semantic gap between the extracted goal term and the operational
meaning of the surrounding Rust --- that gap is
\cref{sec:risks-fidelity}, and full closure means MIR-level
extraction (RefinedRust's move), out of scope for the first campaign
and identified here as the principal residual weakness.

One genuinely new risk arrives with the ghost-token layer: the macro
expands inside the crate under verification, so a bug in
\code{ves-macros} could in principle alter behavior. The mitigations
are the token-identity snapshot tests, the ZST/\code{inline(always)}
discipline (extended into the existing pinned layout-size test
matrices: every token asserts \code{size\_of == 0}), and the fact that
the expansion is deliberately minimal --- \code{let} bindings of unit
structs plus the verbatim subject. This risk class is shared with
every proc macro the workspace already uses, including
\code{dotnet-macros} and \code{dotnetdll}'s own~\cite{dotnetdll}.

\subsection{The model-fidelity gap}
\label{sec:risks-fidelity}

Every theorem is about $\mathcal{M}_{\mathrm{VES}}$,
$\mathcal{M}_{\mathrm{Rust}}$, and hand-stated correspondences --- not
about the compiled binary. The gap has three layers, with different
mitigations:

\begin{itemize}[leftmargin=1.4em]
\item \emph{VES model vs.\ ECMA-335 prose}: mitigated the JSCert
  way~\cite{jscert2014} --- the model is executable and differentially
  tested against real .NET via the existing 486-fixture corpus and diff
  harness, which today has exactly one differential test and deserves
  expansion regardless of this project.
\item \emph{Rust axiomatization vs.\ rustc reality}: the
  $\mathcal{M}_{\mathrm{Rust}}$ axioms are aligned with MiniRust's UB
  list~\cite{minirust} and Tree Borrows~\cite{treeborrows2025}, and
  every axiom gets a Miri-executable counterexample template: if the
  axiom is wrong in a way that matters, a fuzz/Miri run should be able
  to exhibit it. Residual risk is real and stated: theorems are
  conditional on the axiomatization ($\sim$30 axioms is the target
  ceiling; RefinedRust-grade foundations remain the upgrade path for
  critical modules).
\item \emph{Script vs.\ code}: largely mechanized by the design of
  \cref{sec:dsl-elab}. Statement drift (the goal or premises changed)
  moves the statement hash and demotes the site to unproved; subject
  drift (the protected expression changed under an unchanged
  statement) moves the subject hash and flags re-audit. The residue
  is drift the hashes cannot see --- changes to \emph{surrounding}
  code that invalidate a premise without touching the site (a new
  caller that skips a bounds check, a moved guard). The contract-token
  mechanism (\cref{sec:dsl}, templates) closes the caller case by
  typed dataflow; the rest is what the cross-repo drift job and the
  falsifier suites exist to catch.
\end{itemize}

\subsection{Refactor churn under an LLM-agent workflow}
\label{sec:risks-churn}

This codebase refactors at a pace conventional verified projects never
face --- 44-step supervised campaigns are routine here. The embedded
form is markedly better under churn than external annotations:
scripts are ordinary Rust syntax, so \code{rustfmt} formats them,
rename-refactors rewrite spliced identifiers along with everything
else, moving a function moves its proof, and deleting code deletes its
obligations rather than orphaning them. What churn still costs is
proof-manifest staleness --- statements whose hashes move faster than
agents re-prove them. Three design choices exist specifically for
this: obligations factored per \emph{family} rather than per site
(renames move instantiations, not theories); the statement/subject
hash split (site motion re-anchors; only meaning changes re-prove);
and the lag-tolerant manifest policy of \cref{sec:design-compart}, so
\code{cargo build} and the fast CI never wait on Lean. The realistic
posture: proof repair becomes a normal part of the finalize step in
the existing refactor convention, and agents do the first pass --- the
Lean ecosystem is currently the strongest available for that loop. The
corollary is a scoping rule: \emph{annotate the stable core (value
representation, GC boundary, atomics), not the churning periphery
(instruction handlers, intrinsics)}, at least until the model layer
has demonstrated its value.

\subsection{The stop-the-world protocol theorem}
\label{sec:risks-concurrency}

The single theorem most of the corpus leans on --- all mutators parked
$\Rightarrow$ foreign allocations live and stable --- ranges over the
coordinator's lock-backed session model, the thread manager's
handshake, panic-unwind guards (\code{ResumeOnPanic},
\code{CommandCompletionGuard} typestate), safepoint-exclusion counters,
and the lease/generation machinery of arena teardown. This is a
genuine concurrent-protocol verification, the kind of thing
rely-guarantee and Iris ghost state were built
for~\cite{zakowski2018,iris2018}. It is weeks-to-months of work for one
person even with the LTS carefully cut down, and there is no shortcut:
without it, the F1 family rests on a trust grant. The mitigation
is sequencing --- the protocol model (not proof) comes early, gets
model-checked exhaustively for small configurations (2--3 threads) to
expose specification errors cheaply, and only then graduates to a
full proof.
The existing lock-order DAG and the documented forbidden inversions are
a head start: they are, in effect, hand-written lemmas about the LTS.

\subsection{Trusted-premise creep}

The two irreducible trust classes (\cref{sec:obligations-untrusted}) ---
spec-licensed UB from unverifiable IL, and FFI ABI agreement --- are
legitimately axiomatic. The failure mode is sociological: \code{trust} is the
path of least resistance, and a registry that grows monotonically
converts the proof layer back into prose with extra steps. The
\code{TrustGrant} mechanism helps --- an unregistered trusted premise
does not compile, so the registry cannot silently lag the code --- but
the ratchet still requires enforcement: CI reports the grant count as a
budget exactly like the existing \code{check\_mt\_cfg\_ceiling.sh},
and each grant must name the cheapest known falsifier (a Miri mode, a
fuzz target, a fixture) so trust is at least monitored. Two specific
entries deserve flagging now: the \emph{non-moving GC} axiom (fine
today, invalidated the day a compacting collector lands --- the grant
should say so) and the \code{TypeIndex} transmute into a foreign
crate's private layout, which is better fixed upstream in
\code{dotnetdll} (a \code{\#[repr(transparent)]} guarantee or a
constructor) than trusted here --- the author owns that crate too.

\subsection{Provenance}

By design, the managed-pointer serialization calls
\code{expose\_provenance} in its 3-word encoding and so launders
provenance through \code{usize}; strict-provenance Miri is documented
infeasible, and no tool in CI checks provenance obligations today.
The Lean-side $\mathcal{M}_{\mathrm{Rust}}$ can and should model
exposed provenance (the angelic-choice semantics MiniRust specifies),
which makes the obligations \emph{statable}; but the axioms here are
the least falsifiable of the set, and the study flags pointer
serialization/deserialization as the module where proof adds the least
confidence per unit effort until upstream tooling improves.

\subsection{Cost model}
\label{sec:risks-cost}

Order-of-magnitude estimates for a solo developer with strong type
theory, working with LLM agents, part-time on this track
(calibrated against: RefinedRust team-years for foundational unsafe
proofs; Hawblitzel--Petrank two-person SMT-automated GC
verification~\cite{hawblitzelpetrank2009}; JSCert/Jinja multi-year
mechanizations; and the observed productivity of agent-assisted Lean):

\begin{center}
\small
\begin{tabular}{@{}llll@{}}
\toprule
Artifact & Effort & Risk & Standalone value \\
\midrule
CI-gate hardening + defect fixes (a),(b) & \textbf{done} & nil & delivered \\
\code{ves-syntax} + vocabulary v0 + codegen & 2--3 weeks & low & the anchor \\
\code{ves-macros} + \code{ves-tokens} + snapshots & 2--4 weeks & low & rustc-checked premises \\
\code{ves-check} + manifests + ratchet & 1--2 weeks & low & tag-std-class tool \\
$\mathcal{M}_{\mathrm{VES}}$ core, executable + diff-tested & 3--6 months & medium & high (first of its kind) \\
STW/arena protocol LTS + model checking & 1--2 months & medium & high (spec-bug hunting) \\
$\mathcal{M}_{\mathrm{Rust}}$ axiomatization & 1--2 months & medium & moderate \\
Protocol theorem (F1 spine) & 1--3 months & high & the keystone \\
Layout faithfulness (F2) + \code{GcDesc} theorem & 1--2 months & medium & high \\
Critical-site campaign ($\sim$40--60 lemmas) & 3--6 months & medium & the stated goal \\
\bottomrule
\end{tabular}
\end{center}

Summed naively: 1--1.5 years of part-time focused work to the full
stated goal, with usable, independently valuable artifacts landing from
day one --- the first row is already delivered. The embedded-DSL front end adds roughly two to four weeks
over the attribute draft and buys rustc-checked premise discipline in
every build; the estimate's dominant uncertainty remains the protocol
theorem and the degree to which agents can be trusted with Iris-style
ghost reasoning versus plain LTS invariants. The roadmap
(\cref{sec:roadmap}) is cut so that a stall there strands nothing
else.
